\PassOptionsToPackage{table}{xcolor}
\documentclass[10pt,aspectratio=169]{beamer}
\newcommand{\LectureNumber}{11}
\input{course-theme.tex}
\title{For and do-while loops}
\subtitle{CSC103 Programming Fundamentals / Lecture 11}
\author{Dr. Muhammad Farhan}
\institute{Associate Professor of Computer Science\\Department of Computer Science\\COMSATS University Islamabad\\Sahiwal Campus}
\date{Fall 2026}
\begin{document}
\begin{frame}\titlepage\end{frame}

\begin{frame}[fragile,t]{The problem for this lecture}
\begin{columns}[T,onlytextwidth]\begin{column}{0.60\textwidth}
Print a 3 by 4 rectangle of \# characters using nested loops. Explain how many times the inner body runs.
\end{column}\begin{column}{0.35\textwidth}\small
Repeat three times for the three rows.
\end{column}\end{columns}
\note{Introduce the target problem before explaining the tools. Revisit it in the guided solution later.\par Syllabus reading: Deitel Ch 8 (as listed). CLO-2.}
\end{frame}

\begin{frame}[fragile,t]{Learning objectives}
\begin{columns}[T,onlytextwidth]\begin{column}{0.60\textwidth}
\begin{itemize}
\item Select for, while or do-while appropriately
\item Explain break and continue
\item Trace nested iteration
\end{itemize}
\end{column}\begin{column}{0.35\textwidth}\small
CLO-2

The for loop\par The do-while loop\par Break and continue\par Nested loops
\end{column}\end{columns}
\note{Explain how each objective supports the opening problem. The final questions check these objectives.\par Syllabus reading: Deitel Ch 8 (as listed). CLO-2.}
\end{frame}

\begin{frame}[fragile,t]{The for loop}
\begin{itemize}
\item A for header groups initialisation, condition and update.
\item The update runs after each normally completed body.
\item A header variable belongs to the loop scope.
\end{itemize}
\vspace{1mm}\begin{block}{Example}
\begin{lstlisting}
for (int i = 0; i < 4; i++) {
  System.out.println(i);
}
\end{lstlisting}
\end{block}
\note{Order: initialisation once, condition, body, update, then condition again. Output: 0,1,2,3 on separate lines.\par Syllabus reading: Deitel Ch 8 (as listed). CLO-2.}
\end{frame}

\begin{frame}[fragile,t]{The do-while loop}
\begin{itemize}
\item do-while checks its condition after the body.
\item The body therefore runs at least once.
\item The statement ends with a semicolon.
\end{itemize}
\vspace{1mm}\begin{block}{Example}
\begin{lstlisting}
int n = 0;
do {
  System.out.println(n);
  n++;
} while (n < 0);
\end{lstlisting}
\end{block}
\note{This prints 0 once despite the false condition. Menus commonly need an initial display before a decision to repeat.\par Syllabus reading: Deitel Ch 8 (as listed). CLO-2.}
\end{frame}

\begin{frame}[fragile,t]{Break and continue}
\begin{itemize}
\item break exits the nearest loop.
\item continue skips the rest of the current iteration.
\item In a for loop, continue proceeds to the update step.
\end{itemize}
\vspace{1mm}\begin{block}{Example}
\begin{lstlisting}
for (int i = 1; i <= 5; i++) {
  if (i == 2) continue;
  if (i == 4) break;
  System.out.println(i);
}
\end{lstlisting}
\end{block}
\note{Output: 1 and 3. A continue in a while loop can accidentally skip a required update.\par Syllabus reading: Deitel Ch 8 (as listed). CLO-2.}
\end{frame}

\begin{frame}[fragile,t]{Nested loops}
\begin{itemize}
\item The inner loop completes for each outer iteration.
\item With independent bounds, body executions multiply.
\item Choose distinct names for separate counters.
\end{itemize}
\vspace{1mm}\begin{block}{Example}
\begin{lstlisting}
for (int row = 1; row <= 2; row++) {
  for (int col = 1; col <= 3; col++) {
    System.out.print("*");
  }
  System.out.println();
}
\end{lstlisting}
\end{block}
\note{Output is two rows of three stars. This prepares students for two-dimensional array traversal.\par Syllabus reading: Deitel Ch 8 (as listed). CLO-2.}
\end{frame}

\begin{frame}[fragile,t]{The for update follows continue}
\begin{itemize}
\item A for loop transfers from continue to its update expression.
\item A while loop transfers from continue directly to its condition.
\item This difference matters when progress depends on a body statement that continue might skip.
\end{itemize}
\vspace{1mm}\begin{block}{Example}
\begin{lstlisting}
for (...; ...; i++) {
  if (skip) continue;
}
The i++ update still runs.
\end{lstlisting}
\end{block}
\note{Explain the mechanism using the displayed example. A for loop transfers from continue to its update expression. A while loop transfers from continue directly to its condition. This difference matters when progress depends on a body statement that continue might skip.\par Syllabus reading: Deitel Ch 8 (as listed). CLO-2.}
\end{frame}

\begin{frame}[fragile,t]{Row boundaries belong to the outer loop}
\begin{itemize}
\item The inner loop prints all characters in one row.
\item A newline after the inner loop ends that row.
\item Putting the newline inside the inner loop would produce one character per line.
\end{itemize}
\vspace{1mm}\begin{block}{Example}
\begin{lstlisting}
Row 1: #### then newline
Row 2: #### then newline
Row 3: #### then newline
\end{lstlisting}
\end{block}
\note{Explain the mechanism using the displayed example. The inner loop prints all characters in one row. A newline after the inner loop ends that row. Putting the newline inside the inner loop would produce one character per line.\par Syllabus reading: Deitel Ch 8 (as listed). CLO-2.}
\end{frame}


\begin{frame}[fragile,t]{Nested loops visit a rectangular grid}
\begin{center}\begin{tikzpicture}
\node[box,text width=13mm,minimum height=11mm] (n0) at (0.0,0.0) {r0c0};
\node[box,text width=13mm,minimum height=11mm] (n1) at (1.8,0.0) {r0c1};
\node[box,text width=13mm,minimum height=11mm] (n2) at (3.6,0.0) {r0c2};
\node[box,text width=13mm,minimum height=11mm] (n3) at (0.0,-1.2) {r1c0};
\node[box,text width=13mm,minimum height=11mm] (n4) at (1.8,-1.2) {r1c1};
\node[box,text width=13mm,minimum height=11mm] (n5) at (3.6,-1.2) {r1c2};
\node[font=\small,text=accent] at (0.0,.85) {column 0};
\node[font=\small,text=accent] at (1.8,.85) {column 1};
\node[font=\small,text=accent] at (3.6,.85) {column 2};
\node[font=\small,text=accent] at (-1.7,0.0) {row 0};
\node[font=\small,text=accent] at (-1.7,-1.2) {row 1};
\end{tikzpicture}\end{center}
\begin{block}{What to notice}The inner loop completes a row before the outer loop moves to the next row.\end{block}
\note{Trace each labelled connection. The inner loop completes a row before the outer loop moves to the next row.}
\end{frame}
\begin{frame}[fragile,t]{Key distinctions}
\small\renewcommand{\arraystretch}{1.5}
\rowcolors{2}{pale}{white}
\begin{tabularx}{\textwidth}{>{\raggedright\arraybackslash}X>{\raggedright\arraybackslash}X>{\raggedright\arraybackslash}X}
\rowcolor{navy}
\color{white}\bfseries Loop & \color{white}\bfseries Condition check & \color{white}\bfseries Minimum body executions\\
while & Before body & 0\\
for & Before body & 0\\
do-while & After body & 1\\
\end{tabularx}
\vspace{3mm}\footnotesize These distinctions determine which operation or control structure fits the problem.
\note{\par Syllabus reading: Deitel Ch 8 (as listed). CLO-2.}
\end{frame}

\begin{frame}[fragile,t]{First example: execution}
\begin{lstlisting}
int sum = 0;
for (int i = 1; i <= 5; i++) {
  if (i == 3) continue;
  sum += i;
}
System.out.println(sum);
\end{lstlisting}
\note{This is the main body from Java\_Examples/Lecture11.java. The source file includes the complete class. Predict the result before the next slide.\par Syllabus reading: Deitel Ch 8 (as listed). CLO-2.}
\end{frame}

\begin{frame}[fragile,t]{First example: result}
\begin{columns}[T,onlytextwidth]\begin{column}{0.60\textwidth}
12\par The added values are 1, 2, 4 and 5.\par continue still allows the for update to run.
\end{column}\begin{column}{0.35\textwidth}\small
Break and continue

Nested loops
\end{column}\end{columns}
\note{Expected output and interpretation: 12
The added values are 1, 2, 4 and 5.
continue still allows the for update to run.\par Syllabus reading: Deitel Ch 8 (as listed). CLO-2.}
\end{frame}

\begin{frame}[fragile,t]{Guided practice}
\begin{columns}[T,onlytextwidth]\begin{column}{0.60\textwidth}
Print a 3 by 4 rectangle of \# characters using nested loops. Explain how many times the inner body runs.
\end{column}\begin{column}{0.35\textwidth}\small
Attempt the solution before continuing.

Use the definitions and examples from this lecture.
\end{column}\end{columns}
\note{Give students 8 minutes to attempt the problem. The following slides provide a complete solution. Original solution guidance: Outer loop: three iterations. Inner loop: four iterations per row. Print a newline after the inner loop. Total inner body executions: 12.\par Syllabus reading: Deitel Ch 8 (as listed). CLO-2.}
\end{frame}

\begin{frame}[fragile,t]{Solution strategy}
\begin{columns}[T,onlytextwidth]\begin{column}{0.60\textwidth}
\begin{itemize}
\item Repeat three times for the three rows.
\item Inside each row, print four \# characters without ending the line.
\item Print one newline after completing the inner loop.
\end{itemize}
\end{column}\begin{column}{0.35\textwidth}\small
The next program uses fixed inputs so its result can be reproduced.
\end{column}\end{columns}
\note{Discuss why each step is needed. State the input assumptions before executing the program.\par Syllabus reading: Deitel Ch 8 (as listed). CLO-2.}
\end{frame}

\begin{frame}[fragile,t]{Complete solution}
\begin{lstlisting}
public class Solution11 {
  public static void main(String[] args) throws Exception {
    for (int row = 0; row < 3; row++) {
      for (int column = 0; column < 4; column++) {
        System.out.print("#");
      }
      System.out.println();
    }
  }
}
\end{lstlisting}
\note{Complete runnable file: Java\_Examples/Solution11.java. Inputs are fixed in the program.  \par Syllabus reading: Deitel Ch 8 (as listed). CLO-2.}
\end{frame}

\begin{frame}[fragile,t]{Execution trace}
\small\renewcommand{\arraystretch}{1.5}
\rowcolors{2}{pale}{white}
\begin{tabularx}{\textwidth}{>{\raggedright\arraybackslash}X>{\raggedright\arraybackslash}X>{\raggedright\arraybackslash}X>{\raggedright\arraybackslash}X}
\rowcolor{navy}
\color{white}\bfseries Outer row & \color{white}\bfseries Inner column values & \color{white}\bfseries Printed content & \color{white}\bfseries Inner executions\\
0 & 0, 1, 2, 3 & \#\#\#\# & 4\\
1 & 0, 1, 2, 3 & \#\#\#\# & 4\\
2 & 0, 1, 2, 3 & \#\#\#\# & 4\\
All rows & 3 rows  x  4 columns & Three lines & 12\\
\end{tabularx}
\vspace{3mm}\footnotesize Follow the changing state in execution order.
\note{Manually derived trace for the complete solution. Repeat three times for the three rows. Inside each row, print four \# characters without ending the line. Print one newline after completing the inner loop.\par Syllabus reading: Deitel Ch 8 (as listed). CLO-2.}
\end{frame}

\begin{frame}[fragile,t]{Solution result}
\begin{columns}[T,onlytextwidth]\begin{column}{0.60\textwidth}
\#\#\#\#\par \#\#\#\#\par \#\#\#\#
\end{column}\begin{column}{0.35\textwidth}\small
Why this result follows

Print one newline after completing the inner loop.
\end{column}\end{columns}
\note{Expected result: \#\#\#\#
\#\#\#\#
\#\#\#\#\par Syllabus reading: Deitel Ch 8 (as listed). CLO-2.}
\end{frame}

\begin{frame}[fragile,t]{A common incorrect approach}
int i=0;\par while(i\textless{}5) \{\par   if(i==2) continue;\par   i++;\par \}
\vspace{1mm}\begin{block}{Example}
\begin{lstlisting}
At i=2, continue skips the update forever. Move the update before any continue path or use a for loop.
\end{lstlisting}
\end{block}
\note{Ask students to predict the failure before discussing the explanation. The right side gives the correction.\par Syllabus reading: Deitel Ch 8 (as listed). CLO-2.}
\end{frame}

\begin{frame}[fragile,t]{Boundary and failure checks}
\begin{columns}[T,onlytextwidth]\begin{column}{0.60\textwidth}
Check the stated input assumptions.

Print a multiplication table for 1 through 5. Use formatted output and nested loops.
\end{column}\begin{column}{0.35\textwidth}\small
Outer loop: three iterations. Inner loop: four iterations per row. Print a newline after the inner loop. Total inner body executions: 12.
\end{column}\end{columns}
\note{Use the specification to derive each expected result. Exercise solution: Outer loop: three iterations. Inner loop: four iterations per row. Print a newline after the inner loop. Total inner body executions: 12.\par Syllabus reading: Deitel Ch 8 (as listed). CLO-2.}
\end{frame}

\begin{frame}[fragile,t]{Check your understanding}
\begin{columns}[T,onlytextwidth]\begin{column}{0.60\textwidth}
Which loop guarantees one body execution? Does break exit every nested loop?
\end{column}\begin{column}{0.35\textwidth}\small
Write a prediction and explain the rule that supports it.
\end{column}\end{columns}
\note{Answer: do-while. An unlabelled break exits only the nearest enclosing loop or switch.\par Syllabus reading: Deitel Ch 8 (as listed). CLO-2.}
\end{frame}

\begin{frame}[fragile,t]{Answer and explanation}
\begin{columns}[T,onlytextwidth]\begin{column}{0.60\textwidth}
do-while. An unlabelled break exits only the nearest enclosing loop or switch.
\end{column}\begin{column}{0.35\textwidth}\small
At i=2, continue skips the update forever. Move the update before any continue path or use a for loop.
\end{column}\end{columns}
\note{Reconnect the answer to the relevant definition rather than treating it as a fact to memorise.\par Syllabus reading: Deitel Ch 8 (as listed). CLO-2.}
\end{frame}


\begin{frame}[fragile,t]{Compare cases and predict outcomes}
\small\renewcommand{\arraystretch}{1.5}\rowcolors{2}{pale}{white}
\begin{tabularx}{\textwidth}{>{\raggedright\arraybackslash}X>{\raggedright\arraybackslash}X>{\raggedright\arraybackslash}X}\rowcolor{navy}
\color{white}\bfseries Outer index & \color{white}\bfseries Inner indexes & \color{white}\bfseries Body executions\\
0 & 0, 1, 2 & 3\\
1 & 0, 1, 2 & 3\\
Total & 2 rows x 3 columns & 6\\
\end{tabularx}\par\vspace{3mm}\begin{exampleblock}{Discuss}Explain each expected result before running the example. Which case would reveal a wrong assumption?\end{exampleblock}
\note{Read the first two columns and invite predictions before discussing the outcome.}
\end{frame}

\begin{frame}[fragile,t]{Apply the idea}
\begin{alertblock}{Independent practice}Print two rows of three stars. Predict what changes if println() is moved inside the inner loop.\end{alertblock}\vspace{3mm}\begin{enumerate}\item Work through the input by hand.\item Write or correct the program.\item Justify the expected result.\end{enumerate}\note{Allow 3–5 minutes, then compare answers in pairs. Print two rows of three stars. Predict what changes if println() is moved inside the inner loop.}
\end{frame}

\begin{frame}[fragile,t]{Solution and reasoning}
\begin{lstlisting}
for (int row = 0; row < 2; row++) {
  for (int col = 0; col < 3; col++) {
    System.out.print("*");
  }
  System.out.println();
}
\end{lstlisting}\vspace{1mm}\small\begin{itemize}
\item The intended output has two lines, each containing ***.
\item The newline belongs after the inner loop.
\item Moving it inside prints one star per line, producing six lines.
\end{itemize}\note{Answer: The intended output has two lines, each containing ***. The newline belongs after the inner loop. Moving it inside prints one star per line, producing six lines.}
\end{frame}

\begin{frame}[fragile,t]{A multiplication table with nested loops}
\begin{itemize}
\item The row counter supplies one factor and the column counter supplies the other.
\item Reset the column loop for each new row.
\item Formatting belongs inside the inner loop; the newline belongs after it.
\end{itemize}
\vspace{3mm}\footnotesize\textcolor{accent}{Source: supplied ISB campus slides. See speaker notes and ISB\_Source\_Map.md.}
\note{Adapted from the supplied ISB campus folder: Lecture{-}11.pptx, slides 14, 15. Adapted explanation and runnable implementation; sample inputs and traces added for this course.}
\end{frame}

\begin{frame}[fragile,t]{Worked problem: A multiplication table with nested loops}
\begin{block}{Task}Print the 1{-}through{-}3 multiplication table using nested for loops.\end{block}
\small Input: Values are fixed in the program for a reproducible demonstration.\par\vspace{3mm}\begin{exampleblock}{Before running}Predict the output and identify the variables that change.\end{exampleblock}
\note{Adapted from the supplied ISB campus folder: Lecture{-}11.pptx, slides 14, 15. Adapted explanation and runnable implementation; sample inputs and traces added for this course.}
\end{frame}

\begin{frame}[fragile,t]{Complete ISB example (1/1)}
\begin{lstlisting}
public class IsbLecture11 {
  public static void main(String[] args) throws Exception {
    for (int row = 1; row <= 3; row++) {
      for (int col = 1; col <= 3; col++) {
        System.out.printf("%3d", row * col);
      }
      System.out.println();
    }
  }

}
\end{lstlisting}
\note{Adapted from the supplied ISB campus folder: Lecture{-}11.pptx, slides 14, 15. Adapted explanation and runnable implementation; sample inputs and traces added for this course. Complete file: ISB\_Java\_Examples/IsbLecture11.java.}
\end{frame}

\begin{frame}[fragile,t]{Worked trace and expected result}
\small\renewcommand{\arraystretch}{1.4}\rowcolors{2}{pale}{white}
\begin{tabularx}{\textwidth}{>{\raggedright\arraybackslash}X>{\raggedright\arraybackslash}X>{\raggedright\arraybackslash}X}\rowcolor{navy}
\color{white}\bfseries Row factor & \color{white}\bfseries Column factors & \color{white}\bfseries Printed products\\
1 & 1, 2, 3 & 1, 2, 3\\
2 & 1, 2, 3 & 2, 4, 6\\
3 & 1, 2, 3 & 3, 6, 9\\
\end{tabularx}\par\vspace{2mm}
\begin{block}{Expected console output}\ttfamily\footnotesize   1  2  3\par   2  4  6\par   3  6  9\end{block}
\note{Adapted from the supplied ISB campus folder: Lecture{-}11.pptx, slides 14, 15. Adapted explanation and runnable implementation; sample inputs and traces added for this course. Trace and expected values independently worked for this edition.}
\end{frame}

\begin{frame}[fragile,t]{Extend the ISB example}
\begin{alertblock}{Practice}How many inner{-}body executions and newline calls occur for a 10{-}by{-}10 table?\end{alertblock}\vspace{3mm}\begin{exampleblock}{Explain your reasoning}100 product{-}printing executions and 10 newline calls. The two loop bounds multiply.\end{exampleblock}
\note{Adapted from the supplied ISB campus folder: Lecture{-}11.pptx, slides 14, 15. Adapted explanation and runnable implementation; sample inputs and traces added for this course. Answer: 100 product{-}printing executions and 10 newline calls. The two loop bounds multiply.}
\end{frame}
\begin{frame}[fragile,t]{Lecture summary}
\begin{columns}[T,onlytextwidth]\begin{column}{0.60\textwidth}
\begin{itemize}
\item for, while or do-while appropriately
\item break and continue
\item nested iteration
\end{itemize}
\end{column}\begin{column}{0.35\textwidth}\small
\begin{itemize}
\item Repeat three times for the three rows.
\item Inside each row, print four \# characters without ending the line.
\item Print one newline after completing the inner loop.
\end{itemize}
\end{column}\end{columns}
\note{Ask students to give one concrete example for the concepts listed. Use the worked solution to connect the topics.\par Syllabus reading: Deitel Ch 8 (as listed). CLO-2.}
\end{frame}

\begin{frame}[fragile,t]{After-class practice}
\begin{columns}[T,onlytextwidth]\begin{column}{0.60\textwidth}
Print a multiplication table for 1 through 5. Use formatted output and nested loops.
\end{column}\begin{column}{0.35\textwidth}\small
Syllabus reading\par Deitel Ch 8 (as listed)

Next lecture: References and string input
\end{column}\end{columns}
\note{Reading labels follow the supplied outline. Check topic headings against the available textbook edition. Require a program and expected results for the practice task.\par Syllabus reading: Deitel Ch 8 (as listed). CLO-2.}
\end{frame}
\end{document}
